\section{Column control}

\subsection{Control column width and vertical alignment}

\begin{Code}[show-output]
\begin{tblr}{\highLight{|Q[l,2cm,b]|Q[r,2cm,t]|Q[2cm,m]|Q[2cm,m]|Q[2cm,h]|Q[2cm,f]|}}
	\hline
	Column set to \texttt{Q[l,2cm,b]} & Column set to \texttt{Q[r,2cm,t]} & Column set to \texttt{Q[2cm,m]} & Column set to \texttt{Q[2cm,m]} with a lot of text that is going on forever to show the difference between \texttt{t} and \texttt{f} alignment & Column set to \texttt{Q[2cm,h]} & Column set to \texttt{Q[2cm,f]}\\
	\hline
\end{tblr}
\end{Code}

Columns are technically all of \snippet{Q} type. A \snippet{Q} column can take optional arguments in any order, such as:

\begin{itemize}
	\item The horizontal alignment\footnote{A \snippet{Q[l]} column is equivalent to a \snippet{l} column, and so on. However, an \snippet{l} column cannot receive additional arguments, whereas a \snippet{Q} column can: \snippet{Q[l,2cm,f]}.} (\snippet{l}, \snippet{c}, \snippet{r}, \snippet{j}). Default value is \snippet{j}.
	\item The column width (in the example above, in \snippet{cm} but any unit of length, see e.g. \url{https://en.wikibooks.org/wiki/LaTeX/Lengths}, is valid). By default, a column will extend as much as possible so that all cells text fit on one line, which might cause the table to overflow.
	\item The vertical alignment. The possible values are: \snippet{h} (\textbf{h}ead, positions the text at the top of the cell), \snippet{f} (\textbf{f}oot, positions the text at the bottom of the cell), \snippet{t} (sets the baseline of the cell at the \textbf{t}op of the text block), \snippet{m} (sets the baseline at the \textbf{m}iddle of the text block), \snippet{b} (sets the baseline at the \textbf{b}ottom of the text block). Default value is \snippet{m}.
\end{itemize}

Other options are available, and we will see some of these below.

Note that \snippet{h} and \snippet{b} options are not equivalent, neither are \snippet{f} and \snippet{t}.


\subsection{Multiline cells}

\begin{Code}[show-output]
\begin{tblr}{lQ[l,m]Q[l,m]}
	Country & Capital(s)             & Comment \\ \hline
	Belize  & Belmopan               & N/A \\
	Bénin   & \highLight{{Porto-Novo \\ Cotonou}} & \highLight{{Official capital \\ De facto administrative capital}}\\
	Bhutan  & Thimphu                & N/A \\
\end{tblr}
\end{Code}

Multiline cells must be delimited by \{\}. They can be vertically centered, by selecting the desired alignment in the colspec.


\subsection{Variable width column}

\begin{Code}[show-output]
\begin{tblr}{
	\highLight{width=0.9\linewidth},
	colspec={ll|\highLight{X}|\highLight{X[2,r]}|},
}
	Furniture & Price & Shop     & Grade \\ \hline
	Bed       & 150   & MOBILI   & 6/10 \\
	Table     & 800   & LUXTABLE & 8/10 \\
	Desk      & 300   & LBC      & 7/10 \\
\end{tblr}
\end{Code}

\begin{warningbox}{Multiple arguments in a tabularray header}
	The example above passes another argument to tabularray (namely, the table width). Consequently, the colspec \textit{must} be specified by using the keyword \snippet{colspec}. Keywords must be separated by a comma \snippet{,}.
	
	There can be no empty lines in the tabularray header. This would cause a compilation error.
\end{warningbox}


\snippet{X} columns are columns of varying width. tabularray will automatically adjust their width so the table spans the whole line width, or whatever value is the header attribute \snippet{width} set to (if the attribute \snippet{width} is not set, it defaults to the line width).

\snippet{X} columns can receive any other argument that \snippet{Q} columns can receive, e.g. \snippet{X[r]} is a right-aligned \snippet{X} column\footnote{Technically, an \snippet{X[n]} column, with n a coefficient, is an alias to a \snippet{Q[co=n]} column.}.

In case there are several \snippet{X} columns, the extra width is equally distributed among them. This distribution can be weighted by passing any number value (a coefficient) as an argument to the \snippet{X} column. In the example above, the fourth column is twice as large as the third. Default coefficient for an \snippet{X} column is 1.

The coefficient in the \snippet{X} column can be any value (possibly decimal). It can also be negative. For instance, \snippet{X[-2.5]} means that the column will expand \textit{at maximum} with coefficient 2.5, but will be smaller if the text in the column does not require that much space and other \snippet{X} columns can fill the remaining space.


\begin{Code}[show-output]
\begin{tblr}{|X|X[-2.5]|}
	\hline
	coeff +1 & coeff -2.5 \\ \hline
\end{tblr}
\end{Code}


\subsection{Merged cells}

\begin{Code}[show-output]
\begin{tblr}{|l|l|l|l|l|}
	\hline
	1 & 2                                & 3 & \highLight{\SetCell[c=2]{}} 4,5  & hidden!\\ \hline
	1 & \highLight{\SetCell[r=2,c=3]{r,m,4cm}} 2,3,4 & 0 & 0                    & 5 \\ \hline
	\highLight{\SetCell[r=2]{f}} 1 &                 &   &                      & 5 \\ \hline
	0 & 2                                & 3 & 4                    & 5 \\ \hline
\end{tblr}
\end{Code}

To merge cells, you have to use the command \snippet{\textbackslash SetCell} and specify the merging within the optional argument. For instance, \snippet{\textbackslash SetCell[r=2]\{\}} says that the cell will span two rows, starting from the current cell. \snippet{\textbackslash SetCell[c=4]\{\}} will span the current cell over 4 columns.

Merged cells must still be specified in the table body with the ampersands \snippet{\&}, even though their contents will be overwritten. Their contents may be empty.

The mandatory argument of \snippet{\textbackslash SetCell[]\{\}} can receive any cell formatting instruction, e.g. a width, horizontal/vertical alignment, etc. For instance, \snippet{\textbackslash SetCell[r=2,c=3]\{r,m,4cm\}} merges the cells on the next 2 rows and 3 columms, and the resulting merged cell will be right aligned, vertically centered, and 4 cm wide.


\vfill

\begin{infobox}{}
	Most tables do not need more functionalities than what you have seen so far. Feel free to browse the rest of the document for more advanced features.
\end{infobox}

\vfill